\subsection*{Фактор-критические графы}
\begin{definition}
Граф \(G\) \emph{фактор-критический}, если для любой вершины \(u\) граф \(G-u\)
имеет совершенное паросочетание. Тогда \(n\) нечётно, \(\dfc(G)=1\) и
\(\alpha'(G-u)=\alpha'(G)=\frac{n-1}2\) для всех \(u\).
Примеры: \(K_1\), нечётные циклы, \(K_{2k+1}\).
\end{definition}
\begin{lemma}[Галлаи]\label{lem:gallai}
Если \(G\) связен и \(\alpha'(G-u)=\alpha'(G)\) для всех \(u\in V\), то \(G\) фактор-критический.
\end{lemma}
\begin{proof}
По формуле Бержа выберем \(S\) с \(\odd(G-S)-|S|=\dfc(G)\) и предположим \(S\ne\varnothing\).
Пусть \(v\in S\), \(S'=S\setminus\{v\}\). Тогда \((G-v)-S'=G-S\) и
\[
\dfc(G-v)\ge\odd(G-S)-|S'|=\dfc(G)+1 .
\]
Подставляя \(\dfc=v(\cdot)-2\alpha'\): \((n-1)-2\alpha'(G-v)\ge n-2\alpha'(G)+1\), т.\,е.
\(\alpha'(G-v)\le\alpha'(G)-1\) — противоречие с условием. Значит, \(S=\varnothing\) и
\(\dfc(G)=\odd(G)\le1\) (граф связен). Совершенного паросочетания нет (его
удаление любой вершины уменьшило бы \(\alpha'\)), поэтому \(\dfc(G)=1\),
\(n\) нечётно и \(\alpha'(G-u)=\alpha'(G)=\frac{n-1}{2}\): у \(G-u\) совершенное паросочетание.
\end{proof}

\subsection*{Множества \(D\), \(A\), \(C\)}
\begin{definition}[Разбиение Галлаи–Эдмондса]
Для графа \(G\) положим:
\begin{itemize}\tightlist
\item \(D(G)\) — вершины \(u\), для которых существует максимальное паросочетание,
  не покрывающее \(u\) («иногда свободные»);
\item \(A(G)\) — вершины вне \(D(G)\), смежные хотя бы с одной вершиной \(D(G)\);
\item \(C(G)=V\setminus(A(G)\cup D(G))\).
\end{itemize}
\end{definition}
Вершины \(A\cup C\) покрыты \emph{любым} максимальным паросочетанием.
\begin{lemma}[Stability lemma]\label{lem:stability}
Если \(a\in A(G)\), то \(D(G-a)=D(G)\), \(A(G-a)=A(G)\setminus\{a\}\), \(C(G-a)=C(G)\)
и \(\alpha'(G-a)=\alpha'(G)-1\).
\end{lemma}
\begin{proof}
Так как \(a\notin D(G)\), каждое максимальное паросочетание \(M\) покрывает \(a\),
и \(M\) без ребра при \(a\) — паросочетание \(G-a\); при этом \(\alpha'(G-a)\ne\alpha'(G)\)
(иначе максимальное паросочетание \(G-a\) было бы максимальным в \(G\) и не покрывало \(a\)).
Значит, \(\alpha'(G-a)=\alpha'(G)-1\). Равенства для \(A\) и \(C\) следуют из
\(D(G-a)=D(G)\) (по определению \(A,C\) через \(D\)). Докажем его.

\(\supseteq\). Пусть \(v\in D(G)\), \(M_v\) — максимальное паросочетание \(G\), не
покрывающее \(v\), \(ax\in M_v\). Тогда \(M_v\setminus\{ax\}\) имеет размер
\(\alpha'(G)-1=\alpha'(G-a)\), т.\,е. максимально в \(G-a\), и не покрывает \(v\).

\(\subseteq\). Пусть \(M'\) — максимальное паросочетание \(G-a\), не покрывающее
\(v\notin D(G)\). Возьмём \(w\in D(G)\), смежную с \(a\), и максимальное паросочетание
\(M_w\) графа \(G\), не покрывающее \(w\); оно покрывает \(v\) (так как \(v\notin D(G)\)).
В \(M_w\mathbin{\triangle}M'\) компонента \(P\), содержащая \(v\), — путь с началом \(v\),
первое ребро из \(M_w\). Вершина \(a\) не покрыта \(M'\), поэтому лежит не более
чем на одном ребре разности и может быть только концом \(P\).
\begin{enumerate}\tightlist
\item[(а)] \(P\) кончается ребром \(M'\): тогда \(M_w\mathbin{\triangle}E(P)\) — паросочетание
  того же размера, не покрывающее \(v\), т.\,е. \(v\in D(G)\) — противоречие.
\item[(б)] \(P\) кончается ребром \(M_w\) в вершине \(a\): тогда
  \((M_w\mathbin{\triangle}E(P))\cup\{aw\}\) — максимальное паросочетание \(G\),
  не покрывающее \(v\) (вершина \(w\) была свободна) — противоречие.
\item[(в)] \(P\) кончается ребром \(M_w\), \(a\notin P\): тогда \(P\) — \(M'\)-дополняющий путь в
  \(G-a\) (последняя вершина не покрыта \(M'\), так как степень 1), и
  \(|M'\mathbin{\triangle}E(P)|=|M'|+1\) — противоречие с максимальностью \(M'\).\qedhere
\end{enumerate}
\end{proof}
\begin{graphillustration}{Три случая в доказательстве $D(G-a)\subseteq D(G)$. Путь $P$ из $M_w\triangle M'$ начинается в $v$; толстые рёбра — $M_w$, штриховые — $M'$. Вершина $w$ свободна в $M_w$, вершина $a$ — в $M'$.}{fig:06-stability}
\begin{tikzpicture}[font=\small]
\begin{scope}
\node[gi vertex] (v) at (0,0) {$v$};
\node[gi vertex] (p1) at (1,0) {};
\node[gi vertex] (p2) at (2,0) {};
\node[gi vertex] (p3) at (3,0) {};
\draw[gi match] (v) -- (p1); \draw[gi alternate] (p1) -- (p2); \draw[gi match] (p2) -- (p3);
\node[gi vertex] (p4) at (4,0) {};
\draw[gi alternate] (p3) -- (p4);
\node[gi note] at (2,-0.8) {(а) конец ребром $M'$};
\end{scope}
\begin{scope}[shift={(5.2,0)}]
\node[gi vertex] (v) at (0,0) {$v$};
\node[gi vertex] (p1) at (1,0) {};
\node[gi vertex] (p2) at (2,0) {};
\node[gi vertex] (a) at (3,0) {$a$};
\node[gi free] (w) at (3,1.1) {$w$};
\draw[gi match] (v) -- (p1); \draw[gi alternate] (p1) -- (p2); \draw[gi match] (p2) -- (a);
\draw[gi edge,densely dotted] (a) -- (w);
\node[gi note] at (1.6,-0.8) {(б) конец $a$: добавим $aw$};
\end{scope}
\begin{scope}[shift={(9.4,0)}]
\node[gi vertex] (v) at (0,0) {$v$};
\node[gi vertex] (p1) at (1,0) {};
\node[gi vertex] (p2) at (2,0) {};
\node[gi vertex] (p3) at (3,0) {};
\draw[gi match] (v) -- (p1); \draw[gi alternate] (p1) -- (p2); \draw[gi match] (p2) -- (p3);
\node[gi note] at (1.5,-0.8) {(в) $M'$-дополняющий путь};
\end{scope}
\end{tikzpicture}
\end{graphillustration}


\subsection*{Теорема о структуре}
\begin{theorem}[Галлаи 1964, Эдмондс 1965]
Пусть \(D_1,\ldots,D_t\) — компоненты связности \(G[D(G)]\), \(A=A(G)\), \(C=C(G)\). Тогда
\begin{enumerate}\tightlist
\item \(G[C]\) имеет совершенное паросочетание;
\item каждый граф \(D_i\) фактор-критический;
\item любое максимальное паросочетание \(M\) состоит из совершенного паросочетания
  \(G[C]\), почти совершенных паросочетаний всех \(D_i\) и рёбер, соединяющих каждую
  вершину \(A\) с вершинами \emph{различных} компонент \(D_i\);
\item \(\dfc(G)=t-|A|\), \(\ \alpha'(G)=\dfrac{n-t+|A|}{2}\).
\end{enumerate}
\end{theorem}
\begin{proof}
Удаляя вершины \(A\) по одной и применяя лемму~\ref{lem:stability}, получаем
\[
D(G-A)=D(G),\quad A(G-A)=\varnothing,\quad C(G-A)=C,\quad\alpha'(G-A)=\alpha'(G)-|A|.
\]
Так как \(A(G-A)=\varnothing\), в \(G-A\) нет рёбер между \(D\) и \(C\); значит,
\(D_1,\ldots,D_t\) и \(G[C]\) — объединение компонент \(G-A\).

1) Вершины \(C=C(G-A)\) покрыты любым максимальным паросочетанием \(G-A\), а
рёбер из \(C\) в \(D\) нет. Значит, его часть внутри \(C\) — совершенное паросочетание \(G[C]\).

2) Для \(u\in D_i\) есть максимальное паросочетание \(M_u\) графа \(G-A\), не
покрывающее \(u\); его часть внутри компоненты \(D_i\) — максимальное паросочетание
\(D_i\), не покрывающее \(u\). Итак, \(\alpha'(D_i-u)=\alpha'(D_i)\) для всех \(u\), и по
лемме Галлаи~\ref{lem:gallai} \(D_i\) фактор-критический.

3) Пусть \(M\) максимально в \(G\), \(M'\) — \(M\) без рёбер, инцидентных \(A\). Тогда
\(|M'|\ge|M|-|A|=\alpha'(G-A)\), т.\,е. \(M'\) максимально в \(G-A\), и все неравенства —
равенства: каждая вершина \(A\) покрыта своим ребром, и ни одно ребро \(M\) не
соединяет две вершины \(A\). По пунктам 1–2 \(M'\) состоит из совершенного
паросочетания \(G[C]\) и почти совершенных паросочетаний \(D_i\) (в фактор-критическом
графе максимальное паросочетание почти совершенно). Свободных вершин в каждой
\(D_i\) ровно одна, а в \(C\) их нет, поэтому рёбра из \(A\) идут к различным \(D_i\).

4) По пункту 3 свободными остаются ровно \(t-|A|\) вершин — по одной в
компонентах \(D_i\), не сопоставленных с \(A\).
\end{proof}
\begin{graphillustration}{Пример разбиения Галлаи–Эдмондса и одно максимальное паросочетание (толстые рёбра). $D$: треугольник и две одиночные вершины ($t=3$), $A=\{a\}$, $C=\{c_1,c_2\}$. $\dfc=t-|A|=2$: свободны $d_5$ и $d_3$. Ребро $ac_1$ (между $A$ и $C$) не входит ни в одно максимальное паросочетание.}{fig:06-structure}
\begin{tikzpicture}[font=\small]
\node[gi vertex] (a) at (0,0) {$a$};
\node[gi vertex] (d1) at (-2.4,1.3) {$d_1$};
\node[gi vertex] (d2) at (-3.4,2.0) {$d_2$};
\node[gi free] (d3) at (-3.4,0.6) {$d_3$};
\node[gi vertex] (d4) at (-2.4,-0.4) {$d_4$};
\node[gi free] (d5) at (-2.4,-1.5) {$d_5$};
\node[gi vertex] (c1) at (2.2,0) {$c_1$};
\node[gi vertex] (c2) at (3.6,0) {$c_2$};
\draw[gi edge] (d1) -- (d3); \draw[gi edge] (a) -- (d1);
\draw[gi match] (d1) -- (d2);
\draw[gi edge] (d2) -- (d3);
\draw[gi match] (a) -- (d4); \draw[gi edge] (a) -- (d5);
\draw[gi edge] (a) -- (c1); \draw[gi match] (c1) -- (c2);
\node[gi group,fit=(d1)(d2)(d3)(d4)(d5),label=above:{$D(G)$}] {};
\node[gi group,fit=(a),label=above:{$A(G)$}] {};
\node[gi group,fit=(c1)(c2),label=above:{$C(G)$}] {};
\end{tikzpicture}
\end{graphillustration}

\begin{cor}
Ребро с концом в \(D(G)\) входит в некоторое максимальное паросочетание;
ребро внутри \(A\) или между \(A\) и \(C\) — ни в одно.
\end{cor}
\begin{tcolorbox}[colback=cbGrayBack,colframe=cbGray!55]
\textbf{Связь с формулой Бержа.} \(S=A(G)\) — «наилучшее» множество в формуле
Бержа: \(\odd(G-A)\ge t\) (все \(D_i\) нечётны), а \(\odd(G-A)-|A|\le\dfc(G)=t-|A|\).
\textbf{Тривиальные случаи:} есть совершенное паросочетание \(\Rightarrow D=A=\varnothing\);
\(G\) фактор-критический \(\Rightarrow D=V\). Разбиение \(D,A,C\) находится
полиномиально (алгоритм Эдмондса).
\end{tcolorbox}
